\documentclass[instructions.tex]{subfiles}
\begin{document}
 
\chapter{Du bon emploi du temps}
 
On ne doit craindre de perdre le temps, parce que le temps est précieux et parce qu'il est irréparable. 

\primo Vous estimeriez comme une chose infiniment précieuse, un diamant avec lequel vous pourriez acheter un royaume. Le temps de la vie est incomparablement plus estimable, puisque avec ce temps vous pouvez vous procurer un royaume éternel. Le temps ne vous est donné que pour vous sauver; il doit donc vous paraître aussi précieux que le salut. Ce temps est le prix du sang de Jésus-Christ; il n'y a pas un seul moment de ce temps qu'il n'ait demandé sur la croix pour vous à son Père. En perdant le temps, on perd donc les jours de salut, et le prix du sang de l'Homme-Dieu; en le perdant, à quoi s'expose-t-on, puisqu'on se met au hasard de tout perdre! 

Le temps est d'un prix infini. Il vaut en un sens autant que Dieu, dit saint Bernard, puisque la possession de Dieu est la récompense du temps bien employé\footnote{Tempus tantùm valut, quantùm Deus....In tempere benè consumpto comparatur Deus. S.Bern}. Mais ce temps si précieux est cours; et nous n'en avons qu'un moment à la fois, et nous ne savons combien Dieu nous en donnera à l'avenir. Combien donc devons-nous ménager des moments si courts, dont la perte serait si funeste!

Ce temps précieux est cependant la chose que l'on perd avec le plus d'indifférence: les uns, dans des visites et des entretiens inutiles, dans des amusemens frivoles et dans le jeu;, les autres, en négligeant leurs devoirs et les occupations de leur état; d'autres, enfin, dans une molle oisiveté\dots

Que ferons-nous aujourd'hui, disent les fainéans ? il faut nous divertir et passer le temps. Ah! insensés, il ne sera que trop tôt passé pour vous, ce temps. Ne savez-vous pas que vous n'avez pas un moment à perdre, et que vous n'avez peut-être pas deux heures à vivre ? Quand vous seriez assuré de vivre long-temps, la vie la plus longue est toujours un temps bien court pour gagner le ciel. Une éternité même de pénitence, dit saint Augustin, ne serait pas trop longue pour acheter, s'il était possible, la possession d'un Dieu, que vous perdez en abusant du temps\dots

\secundo S'il n'est point de perte plus grave que celle du temps, il n'en est point aussi qui doive plus nous affliger dans l'éternité, parce qu'elle est la plus irréparable. Il n'est point de regret plus sensible dans l'enfer, que d'avoir perdu le temps. Le réprouvé verse des larmes amères, en se souvenant qu'il a perdu tant de momens et d'occasions de se sauver, sans que jamais il puisse obtenir ou réparer un seul des moments perdus. Ah! si un seul instant était encore en son pouvoir pour faire pénitence, quel usage en ferait-il ? Oh! que nous ressentirons cruellement, lorsqu'il n'y aura plus de tempour nous, la perte que nous en faisons à présent ! 

Par un moment de temps on peut acheter le ciel, et mériter de posséder Dieu pour toujours; mais avec toute l'éternité on ne saurait racheter un moment du temps perdu. Un roi puissant, pour prolonger sa vie, donnait jusqu'à cinq cents écus par jour à son médecin. L'insensé ne savait-il pas que le temps ne s'achète point\footnote{Non poterit pretio vel breve tempus enim}. Bientôt le temps sera fini pour vous: vous en comprendrez alors la valeur et la perte, mais trop tard.


\section{Résolutions}

\primo Je distribuerai mon temps de telle sorte, que je puisse en offrir à Dieu tous les instans déjà écoulés. 

\secundo Dans l'emploi que j'en ferai, j'aurai toujours soin de préférer les devoirs de mon état à toute autre chose. 

Mon Dieu,je vous consacre le reste de ma vie; faites que tout ce qui s'en écoulera encore de momens, serve à votre gloire et à mon salut.

\end{document}
